\documentclass[12pt,a4paper]{article}
\usepackage{estilos}
\usepackage{tablas}
\usepackage{estructura}
\usepackage{amsmath}

\newcommand{\tituloTesis}{Implementación del Last Planner System y confiabilidad de la planificación semanal en el mantenimiento vial Limabamba-Challhuayacu, Amazonas}
\newcommand{\campoPortada}{\rule{6cm}{0.35pt}}

\begin{document}
\begin{titlepage}
\begingroup\setstretch{1.0}\setlength{\parskip}{0pt}\centering
{\large\bfseries UNIVERSIDAD PRIVADA ANTENOR\\ORREGO\par}
\vspace{.6cm}{\large ESCUELA DE POSGRADO\par}
\vspace{.45cm}\includegraphics[height=4.7cm]{figuras/escudo.jpeg}\par\vspace{.6cm}
\gradoPortada{PROYECTO DE TESIS PARA OPTAR EL GRADO DE MAESTRO}
\bloqueTituloPortada{\tituloTesis}
\begin{flushleft}
\hspace*{.4cm}Programa: \campoPortada\\[5pt]
\hspace*{.4cm}Línea de investigación: \campoPortada\\[5pt]
\hspace*{.4cm}Sublínea de investigación: \campoPortada\\[14pt]
\hspace*{1cm}\textbf{Autor:} \campoPortada\\[12pt]
\hspace*{1cm}\textbf{Asesor:} \campoPortada\\[12pt]
\hspace*{1cm}Código ORCID: \campoPortada
\end{flushleft}
\vspace{.3cm}\textbf{TRUJILLO, PERÚ}\\[3pt]\rule{1.7cm}{0.35pt}\par
\vspace{.2cm}\begin{flushleft}\hspace*{1.27cm}Modalidad: \campoPortada\end{flushleft}
\endgroup
\end{titlepage}

\setcounter{page}{2}
\encabezadoGeneralidades
\begingroup\setlength{\parskip}{4pt}
\itemGeneral{1.}{Título}
\tituloTesis.
\itemGeneral{2.}{Equipo investigador}
\subitemGeneral{2.1.}{Autor(a)\enspace \campoPortada}
\subitemGeneral{2.2.}{Asesor(a)\enspace \campoPortada}
\itemGeneral{3.}{Tipo de investigación}
Investigación aplicada, de enfoque cuantitativo, alcance correlacional y diseño no experimental, longitudinal y observacional (Hernández Sampieri et al., 2014).
\itemGeneral{4.}{Área y línea de investigación}
\subitemGeneral{4.1.}{Línea\enspace \campoPortada}
\subitemGeneral{4.2.}{Sublínea\enspace \campoPortada}
\itemGeneral{5.}{Unidad académica}
Escuela de Posgrado de la Universidad Privada Antenor Orrego.
\itemGeneral{6.}{Institución y localidad}
Camino vecinal entre los desvíos de Limabamba y Challhuayacu de la ruta AM-676, distrito de Limabamba, provincia de Rodríguez de Mendoza, Amazonas. El expediente técnico identifica al Instituto Vial Provincial de Rodríguez de Mendoza como supervisor del servicio (apartado 1.01).
\itemGeneral{7.}{Periodo de observación}
La observación abarcará las semanas efectivas de ejecución que se identifiquen en los registros del servicio.
% El cronograma académico y el presupuesto se incorporarán con la información indicada en PENDIENTES_VALIDACION.md.
\itemGeneral{9.}{Recursos}
El análisis requerirá acceso autorizado a los registros del servicio, fichas de extracción, un equipo de cómputo y el programa R.
\endgroup
\clearpage
\encabezadoPlan

\rotuloPlan{Resumen ejecutivo del proyecto de tesis}
El mantenimiento periódico del camino vecinal Limabamba-Challhuayacu, en Amazonas, constituye el caso de estudio para determinar la relación entre el grado de implementación del Last Planner System (LPS) y la confiabilidad de la planificación semanal. La investigación será aplicada, de enfoque cuantitativo, alcance correlacional y diseño no experimental, longitudinal y observacional. La unidad de análisis será la semana efectiva de ejecución con registros de planificación y cumplimiento que puedan vincularse. El grado de implementación se evaluará mediante una ficha de análisis documental que abarcará la planificación anticipada, la gestión de restricciones, la planificación semanal colaborativa y el seguimiento de los incumplimientos. La confiabilidad se examinará mediante el Porcentaje de Plan Completado (PPC), calculado como la proporción de compromisos originales terminados íntegramente durante la semana, según los partes de ejecución. Como medidas complementarias se examinarán la estabilidad de los compromisos y la recurrencia de las causas de incumplimiento. El procedimiento seguirá el marco de planificación y control de Ballard y Tommelein (2021). El análisis estimará la magnitud y la dirección de las asociaciones entre las variables a lo largo del periodo de ejecución. Los resultados aportarán evidencia para orientar la revisión de las prácticas de coordinación y control del servicio vial.

\tituloPlan{Problema}
\subtituloPlan{Realidad problemática}
En los proyectos de construcción, la programación general fija la secuencia prevista, mientras que la planificación de corto plazo traduce esa secuencia en compromisos que los responsables pueden asumir. Ballard y Tommelein (2021) sitúan esta tarea dentro del Last Planner System, que articula la preparación anticipada del trabajo, la gestión de restricciones, los compromisos semanales y la revisión de sus incumplimientos.

La calidad de esa revisión depende de cómo se documenta lo ocurrido en obra. En un estudio de campo, Toledo y Sánchez (2024) estudiaron el registro de causas de incumplimiento en proyectos de edificación y propusieron respaldar su clasificación con observaciones de campo. Este antecedente muestra la importancia de vincular las causas reportadas con evidencia de ejecución cuando se examina la repetición de los incumplimientos.

En el Perú, Dileo Diaz et al. (2024) abordaron el empleo del LPS en el mejoramiento de la carretera Oyón-Ambo, tramo I, integrándolo con prácticas de dirección de proyectos y de control económico. Ese antecedente sitúa el sistema en la gestión de infraestructura vial, donde también interesa examinar cómo se sostienen los compromisos de corto plazo en servicios de conservación.

En Limabamba, el expediente técnico del mantenimiento periódico del camino vecinal Limabamba-Challhuayacu describe el alcance técnico y el plazo programado del servicio. La pregunta de investigación se ubica en la ejecución semanal, porque allí se puede comparar qué prácticas de planificación se aplicaron, qué trabajos se comprometieron y cuáles se completaron. Queda por determinar cómo se relaciona el grado de implementación del LPS con la confiabilidad de esos compromisos durante la ejecución del servicio.

\apartadoPlan{Características problemáticas}
El mantenimiento periódico del camino vecinal Limabamba-Challhuayacu comprende una longitud de 11,042 km. El resumen ejecutivo del expediente técnico fija un plazo de 90 días calendario en el apartado 1.09 y describe la reposición del afirmado y los trabajos de drenaje en el apartado 1.06. Estas actividades delimitan el alcance del servicio sobre el cual se examinarán los compromisos semanales.

La confiabilidad se evaluará a partir de los compromisos registrados antes de cada semana y de la evidencia de su terminación, siguiendo el enfoque de Ballard y Tommelein (2021). La gestión de restricciones se observará mediante la asignación de responsables y el seguimiento de requisitos, mientras que la estabilidad se examinará mediante los cambios a los compromisos originales. Para estudiar la recurrencia se vinculará la revisión de los incumplimientos con la reaparición posterior de sus causas. La semana efectiva de ejecución será la referencia común para ordenar estos registros y comparar la planificación con lo ejecutado.

\subtituloPlan{Enunciado del problema}
\apartadoPlan{Problema general}
¿Qué relación existe entre el grado de implementación del Last Planner System y la confiabilidad de la planificación semanal durante el mantenimiento periódico del camino vecinal Limabamba-Challhuayacu, Amazonas?

\apartadoPlan{Problemas específicos}
\begin{listaTrasTitulo}
\item ¿Qué relación existe entre la planificación anticipada mediante \emph{lookahead} y el cumplimiento de los compromisos semanales?
\item ¿Qué relación existe entre la gestión de restricciones y la estabilidad de la planificación semanal?
\item ¿Qué relación existe entre la planificación semanal colaborativa y el Porcentaje de Plan Completado (PPC)?
\item ¿Qué relación existe entre el seguimiento y aprendizaje de los incumplimientos y la recurrencia posterior de sus causas?
\end{listaTrasTitulo}

\tituloPlan{Justificación}
\subtituloPlan{Conveniencia}
La investigación permitirá examinar la planificación del servicio vial a partir de los compromisos que organizan el trabajo de cada semana. El Last Planner System vincula la preparación del trabajo con la programación semanal y la revisión del cumplimiento (Ballard \& Tommelein, 2021). Aplicar este marco al caso facilitará que el equipo responsable consulte, para un mismo periodo, cómo se prepararon las actividades y en qué medida se ejecutaron según lo acordado.

\subtituloPlan{Relevancia social}
Los usuarios del camino vecinal son beneficiarios indirectos del estudio. El Banco Mundial (2010) documentó que el estado de los caminos rurales en el Perú condiciona el acceso a mercados y servicios de salud y educación. La contribución del estudio consiste en aportar información a quienes programan y coordinan el mantenimiento de la vía, de modo que puedan revisar la organización de las actividades previstas para conservarla.

\subtituloPlan{Justificación teórica}
El estudio aportará evidencia sobre la asociación entre las prácticas del LPS y la confiabilidad de la planificación semanal en un servicio de mantenimiento vial rural. A partir del marco de Ballard y Tommelein (2021), se examinará cómo se relacionan la preparación del trabajo, la gestión de restricciones, la coordinación de compromisos y el seguimiento de incumplimientos con los resultados semanales correspondientes a cada práctica. El aporte se circunscribe al servicio estudiado.

\subtituloPlan{Justificación metodológica}
Se establecerá un procedimiento para vincular cada compromiso del plan original con el parte de ejecución que permita verificar su cierre. La ficha de implementación registrará evidencias de las prácticas y la matriz de resultados conservará los compromisos terminados, los cambios y las causas de incumplimiento. La trazabilidad de estos registros permitirá revisar el cálculo semanal del PPC, con el total de compromisos originales como denominador, conforme al criterio de compromisos completados que describen Ballard y Tommelein (2021).

\subtituloPlan{Justificación práctica}
Los resultados podrán orientar la revisión de las reuniones de planificación, del seguimiento de restricciones y de los acuerdos de obra. El registro de causas permitirá examinar los incumplimientos recurrentes mediante evidencia de ejecución, cuya documentación es relevante para clasificarlos (Toledo \& Sánchez, 2024). A partir de las asociaciones observadas, los responsables podrán priorizar la revisión de prácticas concretas de coordinación, asignación de responsables y seguimiento de los compromisos del servicio.

\tituloPlan{Objetivos}
\subtituloPlan{Objetivo general}
Determinar la relación entre el grado de implementación del Last Planner System y la confiabilidad de la planificación semanal durante el mantenimiento periódico del camino vecinal Limabamba-Challhuayacu, Amazonas.

\subtituloPlan{Objetivos específicos}
\begin{listaTrasTitulo}
\item Determinar la relación entre la planificación anticipada mediante \emph{lookahead} y el cumplimiento de los compromisos semanales.
\item Determinar la relación entre la gestión de restricciones y la estabilidad de la planificación semanal.
\item Determinar la relación entre la planificación semanal colaborativa y el Porcentaje de Plan Completado (PPC).
\item Determinar la relación entre el seguimiento y aprendizaje de los incumplimientos y la recurrencia posterior de sus causas.
\end{listaTrasTitulo}


\tituloPlan{Marco teórico}
\subtituloPlan{Antecedentes}
\apartadoPlan{Antecedentes internacionales}
Gallego Bermúdez (2024) estudió la implementación del Last Planner System en la ejecución de una reforma de edificación en Barcelona. El caso sirve como referencia para examinar la aplicación del sistema durante la ejecución de una intervención delimitada.

Brenes Montero (2023) aplicó el LPS durante cuatro semanas en la remodelación de PALI Cañas. Registró restricciones, causas de incumplimiento y la proporción de actividades previstas que se completaron íntegramente. Dos semanas superaron el 80 \% esperado y la cuarta registró 70 \%. El procedimiento orienta el seguimiento semanal de los compromisos y la identificación de sus incumplimientos.

Campos Zúñiga (2024) ensayó el LPS en el proyecto Bodega Z Bridgestone. Comparó las actividades terminadas con las planificadas y registró las causas de incumplimiento. El cumplimiento pasó de 77 \% en la cuarta semana a 100 \% en la quinta y descendió a 88 \% en la sexta. La serie permite discutir la variación del PPC entre semanas consecutivas.

Govindasamy y Bekker (2024) analizaron dos proyectos de infraestructura minera mediante pilotos de LPS. Agregaron mensualmente el PPC semanal y lo compararon con índices de valor ganado. Los PPC medios del piloto fueron 67 \% y 63 \%. La asociación con el índice de desempeño del cronograma presentó valores p de .335 y .462. El estudio permite contrastar el PPC con el desempeño temporal.

Lagos et al. (2024) analizaron 164 proyectos mediante aprendizaje automático. Emplearon indicadores del LPS para predecir el desempeño del cronograma a través de nueve intervalos de ejecución. El modelo seleccionado registró un error absoluto medio de 1,24 \% y un R\textsuperscript{2} medio de .96. El estudio permite discutir cómo los registros sucesivos del LPS informan el seguimiento del proyecto.

\apartadoPlan{Antecedentes nacionales}
Escobedo Cruz (2025) evaluó la aplicación del LPS en partidas críticas de movimiento de tierras y pavimentación del proyecto Ayo-Andamayo mediante un diseño descriptivo. Entre las semanas 28 y 42 reportó un PPC promedio de 85 \% y un porcentaje promedio de restricciones liberadas de 79,65 \%. El caso aporta una referencia vial para discutir ambos indicadores de seguimiento.

Gordillo Cruz y Navio Mejia (2023) evaluaron el LPS y la productividad en el mejoramiento del camino rural Los Valdivia mediante un estudio cuantitativo. Registraron el PPC semanal y clasificaron el tiempo de trabajo por tipo. El PPC pasó de 83,33 \% en la primera semana a 100 \% en la cuarta. El caso aporta un referente para interpretar el cumplimiento en un camino vecinal.

Villacriz Condor (2023) aplicó el LPS en la ampliación de los sistemas de agua y desagüe de Aplao y Huancarqui, Arequipa. Elaboró planes anticipados y semanales y evaluó la confiabilidad de la planificación mediante el PPC y las causas de incumplimiento. El caso orienta el seguimiento de compromisos por frente de trabajo.

Hurtado Manrique (2024) aplicó el LPS en la construcción del parque eólico San Juan de Marcona. Elaboró la programación intermedia, asignó responsables para resolver las restricciones y calculó el PPC a partir de los planes semanales. El caso permite discutir cómo se documenta la preparación de actividades antes de incorporarlas al plan semanal.

Cáceres Jordán (2023) aplicó el LPS en el proyecto rural de riego Ichutorcco mediante un diseño cuasiexperimental. Registró restricciones y calculó el PPC con las actividades programadas y ejecutadas. El PPC fue 68 \% en la primera semana y 54 \% en la segunda. Entre la tercera y la vigésima semana, el estudio reportó valores de 83 \% a 100 \%. La serie permite comparar fluctuaciones del cumplimiento en obras rurales.


\subtituloPlan{Marco conceptual}
\apartadoPlan{Grado de implementación del Last Planner System}
El Last Planner System (LPS) es un sistema de planificación y control de la producción que articula la preparación del trabajo, la formulación de compromisos y el aprendizaje a partir de su ejecución. Su grado de implementación expresa en qué medida se aplican sus prácticas (Ballard \& Tommelein, 2021; Lagos et al., 2019). En esta tesis se evaluará semanalmente mediante evidencias documentales.

\apartadoPlan{Confiabilidad de la planificación semanal}
Es la capacidad del equipo para cumplir los compromisos asumidos en el plan semanal, de modo que las cuadrillas dispongan del trabajo previsto al momento de iniciar sus actividades (Ballard \& Tommelein, 2021). En esta tesis se evaluará principalmente mediante el Porcentaje de Plan Completado (PPC), complementado por el análisis de cambios a los compromisos y de causas recurrentes de incumplimiento.

\apartadoPlan{Planificación anticipada (\emph{lookahead})}
Es la revisión de las actividades previstas para las próximas semanas con el fin de precisar las operaciones, identificar sus requisitos y resolver las restricciones antes de comprometer su ejecución. Vincula la programación de fases con la planificación semanal al preparar actividades que puedan incorporarse al plan de trabajo (Ballard \& Tommelein, 2021; Hamzeh et al., 2012).

\apartadoPlan{Gestión de restricciones}
Comprende identificar las condiciones que impiden iniciar o completar una actividad según lo previsto, asignar responsables y plazos para resolverlas, y dar seguimiento a su liberación. En el LPS, esta práctica permite preparar el trabajo antes de incorporarlo al plan semanal (Ballard \& Tommelein, 2021; Lagos et al., 2019).

\apartadoPlan{Planificación semanal colaborativa}
Es el proceso en el que los responsables directos de ejecutar el trabajo coordinan las actividades de la próxima semana, revisan sus dependencias y acuerdan compromisos conforme a la capacidad y a las condiciones de ejecución. Los acuerdos se incorporan al plan semanal para verificar su cumplimiento al cierre del periodo (Ballard \& Tommelein, 2021).

\apartadoPlan{Seguimiento y aprendizaje de los incumplimientos}
Es la práctica de revisar los compromisos que quedaron sin completar, analizar las causas y acordar acciones para evitar que vuelvan a ocurrir. El aprendizaje se concreta al revisar si las acciones adoptadas modifican las condiciones de ejecución que dieron lugar a los incumplimientos (Ballard \& Tommelein, 2021).

\apartadoPlan{Estabilidad de los compromisos semanales}
El plan semanal recoge los acuerdos de ejecución y sirve de referencia para verificar el cumplimiento (Ballard \& Tommelein, 2021). En esta tesis, la estabilidad se entiende como el grado en que se mantienen los compromisos del plan original durante la semana. Se examinará mediante las modificaciones registradas después del cierre del plan.

\apartadoPlan{Recurrencia de las causas de incumplimiento}
Es la reaparición de una categoría de causa asociada a compromisos incumplidos en distintos periodos de planificación. El registro por categorías permite reconocer problemas repetidos y orientar la búsqueda de sus causas de fondo (Ballard \& Tommelein, 2021). En esta tesis se examinará su reaparición en semanas posteriores a la revisión del incumplimiento.

\subtituloPlan{Términos básicos}
\apartadoPlan{Porcentaje de Plan Completado (PPC)}
Porcentaje obtenido al dividir los compromisos originales terminados íntegramente durante la semana entre el total de compromisos asumidos para ese periodo y multiplicar por cien (Ballard \& Tommelein, 2021).

\apartadoPlan{Restricción}
Condición pendiente, como la disponibilidad de materiales o la conclusión de un trabajo previo, que debe resolverse para ejecutar una actividad según lo previsto (Ballard \& Tommelein, 2021).

\apartadoPlan{Compromiso semanal}
Acuerdo por el que un responsable se obliga a terminar una tarea definida en la fecha prevista en el plan semanal (Ballard \& Tommelein, 2021).

\apartadoPlan{Causa de incumplimiento}
Motivo documentado por el cual un compromiso del plan semanal quedó sin completarse conforme a lo acordado (Ballard \& Tommelein, 2021).


\tituloPlan{Hipótesis}
\subtituloPlan{Hipótesis general}
Existe una asociación entre el grado de implementación del Last Planner System y la confiabilidad de la planificación semanal durante el mantenimiento periódico del camino vecinal Limabamba-Challhuayacu, Amazonas.

\subtituloPlan{Hipótesis específicas}
\begin{listaTrasTitulo}
\item Existe una asociación entre la planificación anticipada mediante \emph{lookahead} y el cumplimiento de los compromisos semanales.
\item Existe una asociación entre la gestión de restricciones y la estabilidad de la planificación semanal.
\item Existe una asociación entre la planificación semanal colaborativa y el Porcentaje de Plan Completado (PPC).
\item Existe una asociación entre el seguimiento y aprendizaje de los incumplimientos y la recurrencia de sus causas en la semana siguiente.
\end{listaTrasTitulo}

\clearpage
\tituloPlan{Variables}
\subtituloPlan{Definición operacional}
El grado de implementación del LPS se medirá como el porcentaje de prácticas verificadas entre las aplicables y evaluables en cada semana. Se calculará para el conjunto de prácticas y para cada dimensión a partir de la ficha de implementación. La matriz de la Tabla~\ref{tab:operacionalizacion} organiza las evidencias documentales por dimensión con base en las prácticas descritas por Ballard y Tommelein (2021).

La confiabilidad de la planificación semanal se medirá principalmente con el PPC. Se examinarán también la proporción de compromisos originales que se mantienen sin cambios y la recurrencia de las causas de incumplimiento en la semana siguiente. El plan original de cada semana fijará el conjunto de compromisos que se seguirá durante ese periodo (Ballard \& Tommelein, 2021).
\[
  \mathrm{PPC}_{t}=100\,\frac{C_t}{P_t},
\]
donde \(P_t\) es el número de compromisos del plan original de la semana \(t\) y \(C_t\) es el número de esos compromisos completados íntegramente durante la misma semana. El resultado se expresa en porcentaje y se calcula para \(P_t>0\) (Ballard \& Tommelein, 2021).

\clearpage
\begin{landscape}\thispagestyle{empty}
\begingroup\setstretch{1.0}\setlength{\parskip}{0pt}
\subtituloPlanConTabla{Matriz de operacionalización de variables}
\refstepcounter{table}\label{tab:operacionalizacion}
\tituloTabla{\thetable}{Operacionalización del grado de implementación del LPS y la confiabilidad de la planificación semanal}
\footnotesize
\iniciotablaHorizontal
\renewcommand{\arraystretch}{1.16}
\noindent\begin{tabularx}{\linewidth}{|
  >{\raggedright\arraybackslash}p{2.55cm}|
  >{\raggedright\arraybackslash}p{3.30cm}|
  >{\raggedright\arraybackslash}p{3.10cm}|
  >{\raggedright\arraybackslash}X|
  >{\raggedright\arraybackslash}p{4.10cm}|
  >{\centering\arraybackslash}p{1.50cm}|}
\hline
\textbf{Variable}&\textbf{Definición operacional}&\textbf{Dimensión}&\textbf{Indicador}&\textbf{Fuentes e instrumentos}&\textbf{Escala}\\\hline
\multirow{4}{2.55cm}{\raggedright Implementación del LPS}&
\multirow{4}{3.30cm}{\raggedright Proporción de prácticas verificadas entre las aplicables y evaluables por semana.}&
Planificación anticipada&
\celdaIndicadores{\item Prácticas verificadas de revisión anticipada de actividades y requisitos.}&
Ficha de implementación. Planes \emph{lookahead}.&Razón (\%)\\\cline{3-6}
&&Gestión de restricciones&
\celdaIndicadores{\item Prácticas verificadas de asignación de responsables y seguimiento de restricciones.}&
Ficha de implementación. Matrices de restricciones y actas de coordinación.&Razón (\%)\\\cline{3-6}
&&Planificación semanal colaborativa&
\celdaIndicadores{\item Prácticas verificadas de coordinación y formulación de compromisos.}&
Ficha de implementación. Actas de coordinación y planes semanales originales.&Razón (\%)\\\cline{3-6}
&&Seguimiento y aprendizaje de los incumplimientos&
\celdaIndicadores{\item Prácticas verificadas de revisión de causas y seguimiento de acciones acordadas.}&
Ficha de implementación. Actas y registros de causas y acciones.&Razón (\%)\\\hline
\multirow{3}{2.55cm}{\raggedright Confiabilidad de la planificación semanal}&
\multirow{3}{3.30cm}{\raggedright Plan original cotejado con ejecución, cambios y causas.}&
Cumplimiento de los compromisos semanales&
\celdaIndicadores{\item PPC: compromisos originales terminados íntegramente respecto de los compromisos originales asumidos.}&
Ficha de compromisos y ejecución. Planes originales y partes de ejecución.&Razón (\%)\\\cline{3-6}
&&Estabilidad de los compromisos semanales&
\celdaIndicadores{\item Compromisos originales mantenidos sin cambios respecto de los compromisos originales asumidos.}&
Ficha de compromisos y ejecución. Versiones del plan y registros de cambios.&Razón (\%)\\\cline{3-6}
&&Recurrencia de las causas de incumplimiento&
\celdaIndicadores{\item Categorías de causas que reaparecen en la semana siguiente respecto de las registradas en la semana de referencia.}&
Ficha de compromisos y ejecución. Registros de causas de semanas consecutivas.&Razón (\%)\\\hline
\end{tabularx}
\fuenteTabla{Elaboración propia}
\paginaHorizontalCentro
\endgroup
\end{landscape}
\clearpage

\tituloPlan{Metodología}
\subtituloPlan{Material}
El expediente técnico delimitará el alcance y el plazo programado del mantenimiento periódico del camino vecinal Limabamba-Challhuayacu, Amazonas. La medición se sustentará en planes anticipados, matrices de restricciones, actas de coordinación, versiones originales de los planes semanales, partes de ejecución y registros de causas de incumplimiento. Se cotejarán las fechas y versiones para identificar los compromisos vigentes al inicio de cada semana y verificar su terminación (Ballard \& Tommelein, 2021).

\subtituloPlan{Población}
La población comprenderá todas las semanas efectivas de ejecución del servicio. El inicio, las suspensiones y la culminación se identificarán mediante actas y partes de obra. El expediente establece un plazo programado de 90 días calendario, equivalente a 12,86 semanas al dividir los 90 días entre siete. La unidad de análisis será la semana efectiva de ejecución.

\subtituloPlan{Muestra}
Se incluirán todas las semanas que dispongan de un plan original fechado antes de su ejecución, evidencia posterior de cumplimiento y registros de las prácticas del LPS que se evaluarán. La cobertura será censal si todas las semanas de la población cumplen estos criterios. Se documentarán las exclusiones por falta de evidencia o inconsistencias irresolubles. Para cada análisis específico se informará el número de semanas con datos para ambos indicadores.

\subtituloPlan{Técnicas e instrumentos de recolección de datos}
Se aplicará análisis documental mediante una ficha de implementación del LPS, una ficha de compromisos y ejecución, y una matriz semanal. La primera registrará prácticas observables de planificación anticipada, gestión de restricciones, planificación colaborativa y seguimiento de incumplimientos. Cada práctica tendrá un criterio de verificación y una fuente identificable. La segunda ficha registrará los compromisos originales, su terminación, las modificaciones y las causas de incumplimiento. El PPC se calculará dividiendo los compromisos originales terminados íntegramente durante la semana entre los compromisos originales asumidos y multiplicando por cien (Ballard \& Tommelein, 2021).

\subtituloPlan{Procedimientos}
Primero, se verificará el periodo de ejecución y se inventariarán los registros de aplicación del LPS. Se preparará el diccionario de datos y se probarán las fichas en un conjunto inicial de registros para uniformar los criterios de extracción. Después, se asignará un identificador a cada compromiso y se vincularán el plan original, sus cambios y los partes de ejecución. Se revisarán duplicados, omisiones y discrepancias. Los datos faltantes se identificarán como tales. Finalmente, se consolidará una fila por semana, con los indicadores y sus fuentes. El PPC quedará sin calcular cuando el plan original carezca de compromisos.

\subtituloPlan{Diseño de contrastación}
La investigación será aplicada, de enfoque cuantitativo, alcance correlacional y diseño no experimental, longitudinal y observacional. Se analizarán mediciones semanales sucesivas de las prácticas del LPS y del cumplimiento de los compromisos a partir de los registros de ejecución (Hernández Sampieri et al., 2014).

\subtituloPlan{Procesamiento y análisis de datos}
El procesamiento se realizará en R. Se describirán los indicadores con medianas, rangos y gráficos semanales. El grado de implementación se estimará mediante la proporción de prácticas verificadas entre las aplicables y evaluables. Se fijarán previamente los criterios de cobertura y se examinará la sensibilidad del puntaje a las diferencias de documentación. Para el contraste general se relacionará dicho puntaje con el PPC, indicador principal de confiabilidad (Ballard \& Tommelein, 2021).

Los contrastes específicos examinarán la planificación anticipada frente al cumplimiento semanal, la gestión de restricciones frente a la proporción de compromisos originales que se mantuvieron sin cambios, y la planificación colaborativa frente al PPC. El seguimiento de incumplimientos se relacionará con la proporción de categorías de causas registradas en una semana que reaparezcan en la siguiente. Se utilizarán pares de semanas consecutivas con registros de causas comparables.

Se utilizará el coeficiente de Spearman para estimar asociaciones monotónicas cuando los indicadores presenten variación suficiente. Se examinarán tendencias y autocorrelación antes de realizar inferencias (Pyper \& Peterman, 1998). Si la serie admite inferencia considerando la dependencia temporal, se empleará un nivel de significancia de 0,05 y se ajustarán los valores p de los contrastes específicos mediante el procedimiento de Holm (1979). En series insuficientes, se reportará el análisis descriptivo y, cuando sea estimable, la asociación observada en el servicio.

\subtituloPlan{Consideraciones éticas}
Se gestionará la autorización de la entidad responsable para consultar los registros internos. Los nombres y demás identificadores personales se sustituirán por códigos en la base de análisis. Los archivos originales se conservarán con acceso restringido y se documentarán las correcciones, exclusiones y decisiones de codificación. Los resultados se presentarán por semana y se protegerá la identidad del personal y la información contractual reservada (ALLEA, 2023).


\tituloPlan{Referencias}
\begingroup\RaggedRight\interlinepenalty=10000
\iniciarReferencias
\referencia{ALLEA. (2023). \emph{The European code of conduct for research integrity} (ed. rev.). \url{https://doi.org/10.26356/ECOC}}
\referencia{Ballard, G., \& Tommelein, I. D. (2021). 2020 current process benchmark for the Last Planner\textregistered{} System of project planning and control. \emph{Lean Construction Journal}, 53–155. \url{https://doi.org/10.60164/47e7h7a1b}}
\referencia{Banco Mundial. (2010). \emph{Peru rural road program: Bringing access to the rural poor}. \url{https://documents1.worldbank.org/curated/en/473581468145191440/pdf/916830BRI0Box30ram0401101000Public0.pdf}}
\referencia{Brenes Montero, F. E. (2023). \emph{Propuesta de mejora en la planificación y control de los proyectos mediante la implementación de la metodología Last Planner para DICOMA Construcción} [Proyecto de graduación, Instituto Tecnológico de Costa Rica]. Repositorio TEC. \url{https://hdl.handle.net/2238/14575}}
\referencia{Cáceres Jordán, Y. M. (2023). \emph{Implementación de la metodología System Last Planner para la mejora de la productividad en el proyecto de riego} [Tesis de licenciatura, Universidad Peruana Unión]. Repositorio Institucional UPeU. \url{https://hdl.handle.net/20.500.12840/7500}}
\referencia{Campos Zúñiga, J. R. (2024). \emph{Implementación de la metodología Last Planner System en los procesos constructivos de la empresa SCH Consultoría y Construcción Tica S.A.} [Proyecto de graduación, Instituto Tecnológico de Costa Rica]. Repositorio TEC. \url{https://hdl.handle.net/2238/15123}}
\referencia{Dileo Diaz, L. E., Altez Galvan, A., Hidalgo Misari, M. A., \& Parillo Coila, V. (2024). \emph{Dirección de proyectos de infraestructura vial integrando la Guía PMBOK\textregistered{}, LPS\textregistered{} y la metodología del resultado operativo “Mejoramiento de la Carretera Oyón Ambo - Tramo I”} [Tesis de maestría, Universidad Peruana de Ciencias Aplicadas]. Repositorio Académico UPC. \url{https://hdl.handle.net/10757/674060}}
\referencia{Escobedo Cruz, S. P. (2025). \emph{Implementación de la metodología Last Planner System para los plazos y costos del proyecto vial Ayo-Andamayo en Arequipa, 2025} [Tesis de pregrado, Universidad Tecnológica del Perú]. Repositorio Institucional UTP. \url{https://hdl.handle.net/20.500.12867/12575}}
\referencia{Gallego Bermúdez, A. F. (2024). \emph{Implementación de la metodología Last Planner System en la ejecución de un proyecto de reforma de edificación. Caso de estudio en la ciudad de Barcelona} [Trabajo de fin de máster, Universitat Politècnica de Catalunya]. UPCommons. \url{https://upcommons.upc.edu/entities/publication/e982dd0e-8d94-4f80-a19a-32b7cd354183}}
\referencia{Gordillo Cruz, E. F., \& Navio Mejia, M. L. (2023). \emph{El Last Planner y la productividad en el mejoramiento del tramo Los Valdivia del camino vecinal rural, Tiabaya, Arequipa, 2023} [Tesis de pregrado, Universidad César Vallejo]. Repositorio Institucional UCV. \url{https://hdl.handle.net/20.500.12692/116893}}
\referencia{Govindasamy, T., \& Bekker, M. C. (2024). Lean construction: Implementing the Last Planner System on mining projects. \emph{Journal of the Southern African Institute of Mining and Metallurgy, 124}(1), 15–24. \url{https://doi.org/10.17159/2411-9717/1695/2024}}
\referencia{Hamzeh, F., Ballard, G., \& Tommelein, I. D. (2012). Rethinking lookahead planning to optimize construction workflow. \emph{Lean Construction Journal}, 15–34. \url{https://doi.org/10.60164/g3c3a1f5c}}
\referencia{Hernández Sampieri, R., Fernández Collado, C., \& Baptista Lucio, M. del P. (2014). \emph{Metodología de la investigación} (6.ª ed.). McGraw-Hill Education.}
\referencia{Holm, S. (1979). A simple sequentially rejective multiple test procedure. \emph{Scandinavian Journal of Statistics, 6}(2), 65–70. \url{https://www.jstor.org/stable/4615733}}
\referencia{Hurtado Manrique, J. C. (2024). \emph{Planificación con Last Planner System de la obra parque eólico San Juan de Marcona} [Tesis de maestría, Universidad Privada Antenor Orrego]. Repositorio Institucional UPAO. \url{https://hdl.handle.net/20.500.12759/22211}}
\referencia{Lagos, C. I., Herrera, R. F., \& Alarcón, L. F. (2019). Assessing the impacts of an IT LPS support system on schedule accomplishment in construction projects. \emph{Journal of Construction Engineering and Management, 145}(10), Article 04019055. \url{https://doi.org/10.1061/(ASCE)CO.1943-7862.0001691}}
\referencia{Lagos, C. I., Herrera, R. F., Mac Cawley, A. F., \& Alarcón, L. F. (2024). Predicting construction schedule performance with last planner system and machine learning. \emph{Automation in Construction, 167}, Article 105716. \url{https://doi.org/10.1016/j.autcon.2024.105716}}
\referencia{Pyper, B. J., \& Peterman, R. M. (1998). Comparison of methods to account for autocorrelation in correlation analyses of fish data. \emph{Canadian Journal of Fisheries and Aquatic Sciences, 55}(9), 2127–2140. \url{https://doi.org/10.1139/f98-104}}
\referencia{Toledo, M. J., \& Sánchez, B. E. (2024). Improving reasons for non-compliance documentation using UAV on construction projects. En D. B. Costa, F. Drevland, \& L. Florez-Perez (Eds.), \emph{Proceedings of the 32nd Annual Conference of the International Group for Lean Construction (IGLC 32)} (pp. 273–282). \url{https://doi.org/10.24928/2024/0106}}
\referencia{Villacriz Condor, B. M. (2023). \emph{Aplicación del sistema del último planificador en una obra de saneamiento en la provincia de Castilla - Arequipa} [Trabajo de suficiencia profesional, Universidad Nacional de Ingeniería]. Repositorio Institucional UNI. \url{https://hdl.handle.net/20.500.14076/27660}}
\endgroup

\end{document}
